\section{Complete formal classification at rational arguments}
\label{sec:classification}

For positive rational $x$, reduce every subscript to lowest terms and set
\begin{equation}\label{eq:eta}
 \eta_x=\xi_{2x}-2\xi_x+\xi_{x/2}.
\end{equation}
By \eqref{eq:J-F}, this is the dilogarithmic element relevant to $J(x)$.
The extra rational multiple of $\pi^2$ has no boundary.

\begin{lemma}[Analytic and formal reciprocity]\label{lem:reciprocity}
For $x>0$,
\begin{equation}\label{eq:J-reciprocity}
 J(x)+J(1/x)=\log^22.
\end{equation}
For positive rational $x$, one also has $\eta_{1/x}=-\eta_x$ in a
common pre-Bloch group. Hence each type of formal reduction is invariant
under $x\mapsto1/x$.
\end{lemma}
\begin{proof}
Integration by parts gives
\[
 J(x)=\log^22-\int_0^1
       \log(1+t)\frac{x t^{x-1}}{1+t^x}\dd t.
\]
The substitution $u=t^x$ identifies the integral with $J(1/x)$. This
known reciprocity also follows from the functional equations discussed
in \cite{CK}.

For the formal statement, \eqref{eq:boundary-F} gives
$\partial\xi_{1/r}=-\partial\xi_r$. Their invariant sum is zero by
\eqref{eq:zero-detection}. Applying this to the three terms of
\eqref{eq:eta} proves the assertion. Thus the formal assertion is not
inferred solely from an equality of numerical periods.
\end{proof}

\begin{theorem}[Complete all-rational classification]\label{thm:classification}
Let $p,q$ be coprime positive integers.

\smallskip\noindent
\textbf{Mixed parity.} Let $e$ denote the even entry and $o$ the odd entry
of $(p,q)$. Then $J(p/q)$ has a formal rational-dilogarithmic reduction
if and only if it has a formal logarithmic reduction, and these are
equivalent to
\begin{equation}\label{eq:mixed-class}
 \boxed{\quad e^2\equiv\pm1\pmod o,
 \qquad o^2\equiv1\pmod{2e}.\quad}
\end{equation}
Congruences modulo $1$ are interpreted as automatic.

\smallskip\noindent
\textbf{Both entries odd.} There is a formal rational-dilogarithmic
reduction exactly when
\begin{equation}\label{eq:odd-class}
 p,q\in\{1,3,5\}.
\end{equation}
There is a formal logarithmic reduction exactly when $p=q=1$.

\smallskip\noindent
Consequently, the rational-dilogarithmic but not logarithmic arguments
are precisely
\begin{equation}\label{eq:six-extras}
 \boxed{\quad\frac13,\ 3,\ \frac15,\ 5,\ \frac35,\ \frac53.\quad}
\end{equation}
\end{theorem}

\begin{proof}
For mixed parity, Lemma~\ref{lem:reciprocity} lets us take $x=e/o$.
Expanding \eqref{eq:boundary-F} gives
\begin{equation}\label{eq:mixed-boundary}
 \partial\eta_{e/o}=D_o(e)-N_e(o),
\end{equation}
where
\begin{align}
 D_o(e)&=\beta_o(2e)-2\beta_o(e)+\beta_o(e/2),\label{eq:D}\\
 N_e(o)&=\beta_{2e}(o)-2\beta_e(o)+\beta_{e/2}(o).\label{eq:N}
\end{align}
The rational wedges cancel because
$-\cl{2e}+2\cl e-\cl{e/2}=0$. The conductors $o$ and $2e$ are
coprime. By Lemma~\ref{lem:norm}, the boundary is rational if and only
if both $D_o(e)$ and $N_e(o)$ vanish. It is then zero, proving that
rational-dilogarithmic and logarithmic reducibility coincide.

Lemma~\ref{lem:second} gives
$D_o(e)=0\Longleftrightarrow e^2\equiv\pm1\pmod o$.
If $N_e(o)=0$, its top-conductor term belongs to
$\Lambda^2\A(\Q(\zeta_e))$. Theorem~\ref{thm:layer} forces
$\beta_{2e}(o)=0$, or $o^2\equiv\pm1\pmod{2e}$.
Since $4\mid2e$ and $o$ is odd, the negative sign is impossible.
Conversely $o^2\equiv1\pmod{2e}$ kills all three symbols of
$N_e(o)$, including the conductor-$1$ or conductor-$2$ degeneracies.
This proves \eqref{eq:mixed-class}.

Now suppose both $p$ and $q$ are odd. Applying
Theorem~\ref{thm:doubling} at both conductors to the boundary expansion
yields
\begin{equation}\label{eq:odd-boundary}
 \begin{split}
 \partial\eta_{p/q}={}&
  \bigl(\beta_q(p)-\beta_q(p/2)\bigr)
 -\bigl(\beta_p(q)-\beta_p(q/2)\bigr)\\
 &+\cl2\wedge\cl{p/q}.
 \end{split}
\end{equation}
For example, the rational contribution before simplification is
$-\cl{2p}\wedge\cl q+2\cl p\wedge\cl q
 -\cl p\wedge\cl{2q}$, which is the last term displayed.
Normalized norm separation and Lemma~\ref{lem:first} show that this
boundary is rational exactly when $p,q\in\{1,3,5\}$.

For logarithmic reduction, normalized norm to $\Q$ first requires
$\cl2\wedge\cl{p/q}=0$. If the reduced odd ratio $p/q$ is not $1$,
the prime-exponent vector of $p/q$ is nonzero and has zero coordinate
at the prime $2$. It is therefore independent of $\cl2$ in
$\A(\Q)$, making this wedge nonzero. If $p=q=1$, all terms vanish
and $J(1)=\tfrac12\log^22$. This proves the odd case and the list
\eqref{eq:six-extras}.
\end{proof}

\subsection{Examples and the missing factor of two}

The mixed pairs $(e,o)=(4,15),(12,5),(12,29)$ satisfy
\eqref{eq:mixed-class}. Thus $J(4/15)$, $J(5/12)$ and $J(12/29)$ have
formal logarithmic reductions. Explicit short logarithmic expressions
for arbitrary pairs of this kind are a separate constructive task;
the neighboring-ratio family will be evaluated below.

By contrast $(e,o)=(8,3)$ satisfies
$8^2\equiv1\pmod3$ and $3^2\equiv1\pmod8$, but fails
$3^2\equiv1\pmod{16}$. The new-layer obstruction can be seen exactly
in $G_{16}$: with $h=9$, the four-term expression
\eqref{eq:four-terms} at $a=3$ is
\[
 (\ee_3-\ee_5)(\ee_1-\ee_7)=2\ee_3-2\ee_5\ne0,
\]
where residues have been identified up to sign. Hence $J(8/3)$ has no
formal rational-dilogarithmic reduction. Testing only modulo $e$ would
produce a false positive.

For $x=2/q$ the theorem recovers exactly the existing classification.
Odd $q$ gives $q=3,5$; a denominator divisible by $4$ reduces to a ratio
with even entry and odd entry $1$; and a denominator $q=2m$ with odd
$m$ gives the odd ratio $1/m$, allowing precisely $m=1,3,5$ for rational
dilogarithms. Thus the positive cases are $q=2,3,5,6,10$ or $4\mid q$;
the cases $6,10$ are not logarithmic in the formal calculus.

\begin{remark}[What has not been proved]\label{rem:periods}
Theorem~\ref{thm:classification} is a necessary and sufficient condition
for a five-term certificate in Definition~\ref{def:formal}. For a ratio
that fails the test, it does not prove that the real number $J(p/q)$ is
linearly independent of all logarithmic products or rational dilogarithm
values. Such a claim would require information about numerical period
evaluation beyond the symbol theorem. In particular, the numerical
transcendence problems in the manuscript remain separate.
\end{remark}
